% Chapter 5: Interests and Geo-Spatial
% Contains: Subject Namespace, Interest Model, Geo-Spatial Features

% ============================================================
\section{Subject Namespace}
\label{sec:subjects}
% ============================================================

Standard subject hierarchy for Avena system and applications:

\begin{Verbatim}[frame=single,fontsize=\small]
avena.                              # System namespace
|-- gossip.                         # State synchronization
|   |-- topology.>                  # Network topology (LoRa-eligible)
|   |-- cost.>                      # Cost metadata (LoRa-eligible)
|   |-- workload.>                  # Workload specs
|   |-- device.>                    # Device registry + positions
|   +-- time.>                      # Time service locations (NTP/PTP)
|-- control.                        # Control plane
|   |-- workload.>                  # Workload commands
|   +-- device.>                    # Device commands
+-- events.                         # System events
    |-- health.>                    # Workload health
    |-- link.>                      # Link up/down
    +-- gc.>                        # Garbage collection

app.                                # Application namespace
+-- <app-name>.                     # Per-application prefix
    +-- ...                         # Application-defined
\end{Verbatim}

\textbf{Default priorities:} \texttt{avena.control.>} $\rightarrow$ HIGH; \texttt{avena.gossip.>} $\rightarrow$ HIGH (small messages); \texttt{avena.events.>} $\rightarrow$ MEDIUM; \texttt{app.>} $\rightarrow$ application-specified via interest.

\textbf{LoRa eligibility:} Subjects under \texttt{avena.gossip.topology.>}, \texttt{avena.gossip.cost.>}, and \texttt{avena.gossip.time.>} are small and infrequent---ideal for LoRa transmission.

% ============================================================
\section{Interest Model}
% ============================================================

\textbf{Core principle:} Data only flows because an interest exists. Publishers emit to subjects; the network routes to subscribers who have registered interest. No interest means no delivery, no storage, no relay.

This is pull-based semantics with push delivery. Subscribers declare what they need and how they need it.

\begin{lstlisting}[language=Rust]
struct Interest {
    id: InterestId,
    owner: DeviceId,

    // What data
    subject: SubjectPattern,

    // Delivery urgency (use lowest that works)
    priority: Priority,

    // Retention semantics
    retention: RetentionPolicy,

    // For "latest with history" pattern
    historic_priority: Option<Priority>,

    // Time bounds
    ttl: Option<Duration>,

    // Ordering
    ordering: Ordering,

    // Geo-spatial bounds (optional)
    geo_bounds: Option<GeoBounds>,

    // Cost routing override (optional)
    cost_profile: Option<CostProfile>,

    // CRDT metadata
    hlc: HLC,
    terminated: bool,
}

enum RetentionPolicy {
    All,                // Every message
    Latest,             // Only most recent per subject
    LatestN(u32),       // Last N messages
    Window(Duration),   // Messages within time window
}

enum Ordering {
    Ordered,    // Deliver in HLC order (buffered at receiver)
    Unordered,  // Deliver as available (faster)
}
\end{lstlisting}

\textbf{Ordered delivery:} When \texttt{Ordering::Ordered} is specified, messages buffer at the receiver and are delivered in HLC order. This introduces latency when messages arrive out of order. An optional \texttt{max\_latency} parameter bounds waiting time; if the expected prior message does not arrive within max\_latency, it is considered lost and delivery proceeds.

\textbf{Subscriber responsibility:} Use the lowest priority that meets requirements. This relieves network burden and allows critical traffic to flow.

\textbf{Multiple interests:} Same subject can have multiple interests with different semantics. One subscriber gets \texttt{Latest} for real-time display; another gets \texttt{All} for historical records.

\textbf{Delivery set:} For interest $i$ on device $d$ at time $t$, the messages to deliver:
\begin{align*}
\text{deliver}(i, d, t) = \{m : &\ \text{subj}(m) \in \text{pattern}(i) \\
&\land \text{permitted}(d, \text{subj}(m)) \\
&\land \text{satisfies\_retention}(m, i) \\
&\land \text{satisfies\_geo}(m, i, d, t) \\
&\land \lnot\text{delivered}(m, i) \\
&\land \lnot\text{terminated}(i)\}
\end{align*}

\subsection{Retention Patterns}

\begin{table}[H]
\centering
\begin{tabular}{@{}llll@{}}
\toprule
\textbf{Pattern} & \textbf{Priority} & \textbf{Retention} & \textbf{Use Case} \\
\midrule
Real-time position & HIGH & Latest & Current machine location \\
Telemetry stream & MEDIUM & All & Sensor readings for analysis \\
Position history & LOW & All & GPS breadcrumbs for records \\
Latest + history & HIGH/LOW & All & Fast current + eventual history \\
Heartbeat & LOW & Latest & Device alive signal \\
Critical alert & HIGH & All & Events that must not be lost \\
V2X BSM & HIGH & Latest & Safety-critical vehicle position \\
SPaT & HIGH & Latest & Signal phase/timing \\
\bottomrule
\end{tabular}
\caption{Common retention patterns}
\end{table}

\textbf{Latest with history pattern:}

\begin{lstlisting}[language=Rust]
Interest {
    subject: "app.position.tractor-1",
    priority: Priority::High,
    retention: RetentionPolicy::All,
    historic_priority: Some(Priority::Low),  // Older messages demoted
}
\end{lstlisting}

When new message arrives: deliver immediately at HIGH priority. Older undelivered messages remain queued but at LOW priority, yielding to other traffic.

% ============================================================
\section{Geo-Spatial Features}
% ============================================================

\subsection{Position Tracking}

Devices report position via gossip (LoRa-eligible for updates):

\begin{lstlisting}[language=Rust]
struct GeoPosition {
    latitude: f64,
    longitude: f64,
    altitude_m: Option<f32>,
    accuracy_m: Option<f32>,
}
\end{lstlisting}

Stationary devices have fixed positions configured at deployment. Mobile devices update position periodically.

\subsection{Geo-Bounded Interests}

Subscribers can limit data flow based on publisher location:

\begin{lstlisting}[language=Rust]
struct GeoBounds {
    boundary: Option<GeoBoundary>,
    relative_radius_m: Option<u32>,  // Relative to subscriber's position
}

enum GeoBoundary {
    Radius { center: GeoPoint, radius_m: u32 },
    Polygon { vertices: Vec<GeoPoint> },
    Named { boundary_id: String },  // e.g., "field-north-40"
}
\end{lstlisting}

\textbf{Use case --- Agricultural traffic optimization:}

\begin{lstlisting}[language=Rust]
// Only receive high-rate telemetry when tractor-2 is within 300 feet
Interest {
    subject: "app.telemetry.tractor-2.high-rate",
    priority: Priority::High,
    retention: RetentionPolicy::Latest,
    geo_bounds: Some(GeoBounds {
        relative_radius_m: Some(100),  // ~300 feet
    }),
}
\end{lstlisting}

\textbf{Use case --- V2X intersection awareness:}

\begin{lstlisting}[language=Rust]
// Receive SPaT from signals within 500 feet
Interest {
    subject: "v2x.spat.>",
    priority: Priority::High,
    retention: RetentionPolicy::Latest,
    geo_bounds: Some(GeoBounds {
        relative_radius_m: Some(150),  // ~500 feet
    }),
}
\end{lstlisting}

\subsection{Geo-Spatial Permissions}

Certificates can include geo-constraints on subject permissions:

\begin{lstlisting}[language=Rust]
struct SubjectPermission {
    pattern: SubjectPattern,
    action: Action,  // Publish, Subscribe
    geo_constraint: Option<GeoConstraint>,
}

struct GeoConstraint {
    publisher_bounds: Option<GeoBoundary>,
    subscriber_bounds: Option<GeoBoundary>,
    mutual_proximity_m: Option<u32>,
}
\end{lstlisting}

\textbf{Position staleness:} Position information propagates via gossip and may be stale relative to message timestamps. Avena adopts receiver-side filtering: messages route optimistically, and the receiving device performs final geo validation using its current position. This trades potential over-delivery for simplicity. Applications requiring strict geo-fencing should implement additional validation.

